\lysh{38816}{4}

\begin{alist}
\item Η Μαρία κέρδισε:
$[3 \cdot (9 - 1) + 1200] + [5 \cdot (7 - 1) + 1200] = (3 \cdot 8 + 1200) + (5 \cdot 6 + 1200) = 1224 + 1230 = 2454$ βαθμούς.

\item
\begin{rlist}
\item Εφόσον $x$ είναι το πλήθος των αντιπάλων του Θρασύβουλου και $t$ η διάρκεια της παρτίδας, σε λεπτά, τότε:
$x(t - 100) + 800 = 400$ ή $x(t - 100) = -400$ ή $t - 100 = -\dfrac{400}{x}$ ή $t = 100 - \dfrac{400}{x}$.

\item Ο χρόνος της παρτίδας είναι θετικός αριθμός, άρα $t > 0$ ή $100 - \dfrac{400}{x} > 0$, άρα η ανισότητα είναι αληθής.

\item Από το ii. έχουμε ότι $100 > \dfrac{400}{x}$ ή $100x > 400$ ή $x > 4$.
Ο $x$ είναι θετικός ακέραιος, καθώς εκφράζει το πλήθος των αντιπάλων που κέρδισε ο Θρασύβουλος, άρα $x \ge 5$.
Οι παίκτες σε μια παρτίδα είναι το πολύ $12$. Άρα ο Θρασύβουλος κέρδισε το πολύ $11$ αντιπάλους. Επομένως $x \le 11$.
Άρα οι αντίπαλοι που κέρδισε ο Θρασύβουλος ήταν από $5$ μέχρι $11$.

\item Αν ο Θρασύβουλος έκανε αυτή τη νίκη στη $10$η φάση του φετινού πρωταθλήματος, τότε θα κέρδιζε $10x + 1200$, όπου $x$ είναι οι ηττημένοι της παρτίδας.
Αν κέρδιζε $1240$ βαθμούς, τότε:
$10x + 1200 = 1240$ ή $10x = 40$ ή $x = 4$, που απορρίπτεται, γιατί σε αυτή την παρτίδα οι αντίπαλοι ήταν από $5$ έως $11$.
\end{rlist}
\end{alist}

